\documentclass[../../../main.tex]{subfiles}
\begin{document}
\subsection{Time Evolution}
\subsubsection{Schrödinger picture.}
This corresponds to active transformation where the  state is being act upon by the unitary transformation $\ket{\psi (t)} =  U(t)\ket{\psi}$.
The observable is fixed $    O (t)= O$ and the expectation value reads 
\begin{equation*}
\langle O\rangle(t) = \langle\psi(t)|O|\psi(t)\rangle
=\braket{\psi | U ^\dagger (t)O U(t)|\psi}
\end{equation*}

The Schrödinger equation can be derived from this view alongside with unitary operator structure
\begin{equation*}
\frac{d}{dt}\ket{\psi(t)} = \frac{d e^{-iHt/\hbar}}{dt}\ket{\psi(0)}
= -\frac{i}{\hbar}H\ket{\psi(t)}
= -\frac{i}{\hbar}H\ket{\psi(t  )}
\end{equation*}


\subsubsection{Heisenberg picture.}
Here, the observables is the unitary operator act upon $O(t) = U ^\dagger (t)O U(t)$.
The state is fixed $\ket{\psi (t)} = \ket{\psi}$, so this is the passive picture and the expectation value remains the same
\begin{equation*}
\langle O\rangle(t) = \langle\psi|O(t)|\psi\rangle
=\braket{\psi | U ^\dagger (t)O U(t)|\psi}
\end{equation*}
Then the Heisenberg equation 
\begin{equation*}
    \frac{dO_H}{dt}     =  \frac{1}{i \hbar} [O(t),H(t)]
\end{equation*}
is derived as 
\begin{align*}
    \frac{dO_H}{dt} 
    & =  \frac{de^{iHt/\hbar}}{dt}\,O\,U + U^\dagger\,O\,\frac{de^{-iHt/\hbar}}{dt}\\
    & =  \frac{i}{\hbar} (U ^\dagger H O U - U ^\dagger O H U)\\
    & =  \frac{i}{\hbar} (U ^\dagger H U U ^\dagger O U - U ^\dagger O U U ^\dagger H U)\\
    \frac{dO_H}{dt} 
    & =  \frac{i}{\hbar} [H(t) O(t) - O(t) H(t)]  =  \frac{1}{i \hbar} [O(t),H(t)]\\
\end{align*}

\subsection{Compatibility}
The compatibility measure the degree of incompatibility between two observables; it is defined by the commutator
\begin{equation*}
    [\Omega\Lambda]=\Omega\Lambda-\Lambda\Omega
\end{equation*}

\subsubsection{Procedure.}
Here are the scheme by which the compatibility of two operators is measured.
Let us first measure $\Omega$ on the ensemble described by $\ket{\psi}$ and take the particles that yield a result $\omega$.
We immediately measure $\Lambda$ and pick those particles that give a result $\lambda$.
\begin{equation*}
    \ket{\psi}\xrightarrow{\Omega=\omega}\ket{\omega}\xrightarrow{\Lambda=\lambda}\ket{\lambda}
\end{equation*}
Here the two operator are said to be incompatible since the second measurement change the state produced by the first measurement.

Opposite case occur if the first eigenstate $\ket{\omega }$ is also an eigenstate of the second operator $\Lambda$.
We denote this simultaneous eigenstates as $\ket{\omega\lambda}$.
\begin{equation*}
    \ket{\psi}\xrightarrow{\Omega=\omega}\ket{\omega\lambda}\xrightarrow{\Lambda=\lambda}\ket{\omega\lambda}
\end{equation*}
This eigenstate obey the following equations
\begin{align*}
    \Omega \ket{\omega\lambda} & = \omega \ket{\omega\lambda}  \\
    \Lambda\ket{\omega\lambda} & = \lambda \ket{\omega\lambda}
\end{align*}
Now operate the first $\Lambda$ and the second by $\Omega$
\begin{align*}
    \Lambda\Omega \ket{\omega\lambda} & = \omega\lambda \ket{\omega\lambda} \\
    \Omega\Lambda\ket{\omega\lambda}  & = \omega\lambda \ket{\omega\lambda}
\end{align*}
then subtract the second with the first
\begin{align*}
    (\Omega\Lambda-\Lambda\Omega)\ket{\omega\lambda} & = (\omega\lambda-\omega\lambda)\ket{\omega\lambda} \\
    [\Omega,\Lambda]\ket{\omega\lambda}              & = 0 \ket{\omega\lambda}
\end{align*}
Thus $[\Omega,\Lambda]$ must have eigenkets with zero eigenvalue if simultaneous eigenkets are to exist.

\subsection{Hamiltonian}
\subsubsection{Constructing Hamiltonian.}
Hamiltonian is constructed by making the substitution $\mathcal{H}(x,p) \rightarrow H(X,P)$, where $\mathcal{H }$ is the classical Hamiltonian for the same case.

Consider a harmonic oscillator.
The classical Hamiltonian is
\begin{equation*}
    \mathcal{H }=\frac{p^2 }{2m }+\frac{1 }{2 }m\omega^2x^2
\end{equation*}
Thus the Hamiltonian for quantum harmonic oscillator is
\begin{equation*}
    H=\frac{P^2 }{2m }+\frac{1 }{2 }m\omega^2X^2
\end{equation*}
or in three dimension
\begin{equation*}
    H=\frac{P_x^2+P_y^2+P_z^2}{2m  }+\frac{1 }{2}m\omega \left[ X^2+Y^2+Z^2 \right]
\end{equation*}
assuming the force constant is the same in all directions.
It is the isotropic harmonic oscillator, meaning the restoring force has the same frequency $\omega$ in every spatial direction.
This is derived from the potential
\begin{equation*}
    V(x,y,z) = \frac{1}{2} m \omega^2 (x^2 + y^2 + z^2)
\end{equation*}
which generate the isotropic force
\begin{equation*}
    \mathbf{F }= -m \omega^2 \mathbf{r}
\end{equation*}

Now consider a particle in one dimension is subject to a constant force $f$.
The classical Hamiltonian is
\begin{equation*}
    \mathcal{H }=\frac{p^2 }{2m }-fx
\end{equation*}
Thus the quantum Hamiltonian
\begin{equation*}
    H=\frac{P^2 }{2m }-fX
\end{equation*}

Now consider particle of charge $q$ in an electromagnetic field.
The classical Hamiltonian is
\begin{equation*}
    \mathcal{H }=\frac{|\mathbf{p }-(q/c) \mathbf{A}(\mathbf{r},t)|^2 }{2m }+q\phi(\mathbf{r },t)
\end{equation*}
The Hamiltonian is constructed by using the symmetric form
\begin{equation*}
    H=\frac{P^2 }{2m }\left( \mathbf{P }\cdot \mathbf{P }-\frac{q }{c }\mathbf{P }\cdot \mathbf{A}-\frac{q }{c}\mathbf{A}\cdot \mathbf{P}+\frac{q^2 }{c^2}\mathbf{A }\cdot \mathbf{A}\cdot \mathbf{A} \right) +q\phi
\end{equation*}

\subsection{Parity}
Parity refers to the transformation of spatial inversion about the origin $x \mapsto -x$.
The corresponding operator is called the parity operator $\Pi$ defined by its action on the position eigenkets
\begin{equation*}
    \Pi \ket{x }=\ket{-x}
\end{equation*}
Given a state $\ket{\psi}$ and its position representation $\psi(x) = \braket{x|\psi}$, the action of the parity operator on this wavefunction is
\begin{equation*}
    \Pi\psi(x) = \psi(-x) .
\end{equation*}

If $\ket{\psi}$ is an eigenstate of the parity operator, then under spatial inversion it transforms as
\begin{equation*}
    \psi(-x) = \lambda   \psi(x) ,
\end{equation*}
where $\lambda$ is the corresponding eigenvalue.
Since $\Pi^2 = I $ its eigenvalues can only be $\lambda = \pm 1$.
This leads to two possibilities: even parity (symmetric) and odd parity (antisymmetric)
\begin{equation*}
    \psi(-x) = +\psi(x) \qquad    \psi(-x) = -\psi(x)
\end{equation*}
In quantum mechanics parity means \enquote{eigenvalue under the inversion symmetry operator}, while mathematically it means \enquote{even} and \enquote{odd} function.

\subsubsection{Symmetric potential.}
If the potential satisfies
\begin{equation*}
    V(-x) = V(x) ,
\end{equation*}
then the Hamiltonian operator
\begin{equation*}
    H = -\frac{\hbar^2}{2m}\frac{d^2}{dx^2} + V(x)
\end{equation*}
commutes with the parity operator:
\begin{equation*}
    [H, \Pi] = 0
\end{equation*}
This can be easily seen  by considering that kinetic term and the potential are invariant.
Then, eigenfunctions of $H$ (stationary states) can be chosen to simultaneously be eigenfunctions of $\Pi$.
Thus, its eigenfunctions are exactly the even and odd function
\begin{equation*}
    \psi(-x) = +\psi(x) \qquad    \psi(-x) = -\psi(x)
\end{equation*}

\subsection{WKB Approximation}
Wentzel--Kramers--Brillouin (WKB) is a technique to obtain approximate solutions to time independent problems, mainly in 1D or where only $r$ matters in 3D.
The WKB method is valid when the potential varies slowly on the scale of the particle’s wavelength.

In classically allowed regions $V(x)<E$, WKB provides an oscillatory wavefunction of the form
\begin{equation*}
    \psi(x) \sim \frac{1}{\sqrt{p(x)}} \left[A\exp\left(\frac{i}{\hbar} \int^{x} p(x') dx'\right) + B\exp\left(-\frac{i}{\hbar} \int^{x} p(x') dx'\right)\right]
\end{equation*}
In classically forbidden regions $E<V(x)$, WKB provides an exponentially decaying or growing wavefunction
\begin{equation*}
    \psi(x) \sim \frac{C}{\sqrt{|p(x)|}} \left[A\exp\left( \frac{1}{\hbar} \int^{x} |p(x')| dx' \right) + B\exp\left( -\frac{1}{\hbar} \int^{x} |p(x')| dx' \right)\right]
\end{equation*}
The momentum is simply the classical momentum
\begin{equation*}
    p_{V<E}(x) = \sqrt{2m[E - V(x)]}\qquad
    p_{E<V}(x) = \sqrt{2m[ V(x)-E]}
\end{equation*}

\subsubsection{WKB quantization.}
The WKB quantization condition depends on the nature of the boundaries: for two smooth turning points
\begin{equation*}
    \int p(x) dx = (n + \frac{1}{2}) \pi \hbar
\end{equation*}
two hard walls (infinite well)
\begin{equation*}
    \int p(x) dx = n \pi \hbar
\end{equation*}
one hard wall and one smooth turning point
\begin{equation*}
    \int p(x) dx = (n + \frac{3}{4}) \pi \hbar
\end{equation*}
Note that a turning point is called smooth if, in a neighborhood of $x_t$, $V(x)$ is finite and differentiable, while $V'(x)$ finite and nonzero.

\subsection{Perturbation Theory}
Time-independent perturbation theory is an approximation scheme that applies in the following context: we know the solution to the eigenvalue problem of the Hamiltonian $H^0$, and we want the solution to
\begin{equation*}
    H=H^0 +\epsilon H^1
\end{equation*}
where $H^1$ is small compared to $H^0$.
Note that the assumption is not that $\epsilon$ is small, rather $H^1$ is small compared to $H^0$, or more accurately
\begin{equation*}
    \left|\frac{\left\langle m^0 | H | n^0 \right\rangle}{E_n^0 - E_m^0} \right| \ll 1
\end{equation*}

We assume that to every eigenket $|E_n^0\rangle \equiv |n^0\rangle$ of $H^0$ with eigenvalue $E_n^0$, there is an eigenket $|n\rangle$ of $H$ with eigenvalue $E_n$.
We then assume that the eigenkets and eigenvalues of $H$ may expanded in a perturbation series:
\begin{align*}
    |n\rangle & =
    |n^{0}\rangle +\epsilon |n^{1}\rangle +\epsilon^2 |n^{2}\rangle + \dots \\
    E_n       & =
    E_n^{0} + \epsilon E_n^{1} +\epsilon^2 E_n^{2} + \dots
\end{align*}
Although not physical states themselves, $\ket{n^k}$ have physical interpretation: $\ket{n^0}$ the unperturbed state;  $\ket{n^2}$ describes first-order mixing with other unperturbed states;  and $\ket{n^2}$ describes second-order mixing effects, including indirect coupling.

The result of zeroth order approximation is our previous assumption
\begin{equation*}
    H^0 \ket{n^0}=E_n^0 \ket{n ^0}
\end{equation*}
The result of first order approximation are
\begin{align*}
    \ket{n}
          & =\ket{n^0}+\sum_{m \neq n} \frac{ \braket{m^{0} | H^{1} | n^{0} }}{E^{0}_n - E^{0}_m}\ket{m^0} \\
    E^1_n & =  \langle n^0 | H^1 | n^0 \rangle
\end{align*}
Using the result of first order approximation, the second order approximation can be obtained as follows
\begin{align*}
    E_n^{2} = \sum_{m \neq n} \frac{|\langle n^0 | H | m^0 \rangle|^2}{E^0_n - E^0_m}
\end{align*}

\subsubsection{Derivation.}
We then substitute these into the eigenvalue equation
\begin{multline*}
    (H^0+\epsilon H^1) \left( |n^{0}\rangle +\epsilon |n^{1}\rangle +\epsilon^2 |n^{2}\rangle + \dots  \right) \\
    = \left( E_n^{0} + \epsilon E_n^{1} +\epsilon^2 E_n^{2} + \dots \right) \left( |n^{0}\rangle +\epsilon |n^{1}\rangle +\epsilon^2 |n^{2}\rangle + \dots  \right)
\end{multline*}
Then collect the same power of $\epsilon$.
The zeroth order comes from  $\epsilon^0$
\begin{equation*}
    H^0 \ket{n^0}=E_n^0 \ket{n ^0}
\end{equation*}
By assumption, this equation may be solved and, so we move on to the first-order terms, which comes from $\epsilon^1$
\begin{equation*}
    H^{0} | n^{1} \rangle + H^{1} | n^{0} \rangle = E^{0}_{n} | n^{1} \rangle + E^{1}_{n} | n^{0} \rangle
\end{equation*}

Dotting both side with $\bra{n^0}$
\begin{align*}
    \langle n^0 | H^1 | n^1 \rangle + \langle n^0 | H^1 | n^0 \rangle & =  E^0_n \langle n^0 | n^1 \rangle + E^1_n \langle n^0 | n^0 \rangle \\
    E^0_n \langle n^0 | n^1 \rangle + \langle n^0 | H^1 | n^0 \rangle & =  E^0_n \langle n^0 | n^1 \rangle + E^1_n                           \\
    E^1_n                                                             & =  \langle n^0 | H^1 | n^0 \rangle
\end{align*}
This result state the first-order change in energy is the expectation value of $H^1$ in the unperturbed state.

Let us next dot both sides with another eigenket of the same unperturbed Hamiltonian $\bra{m^0}$, with $m \neq n$
\begin{align*}
    \langle m^{0}|H^{0}|n^{1}\rangle + \langle m^{0}|H^{1}|n^{0}\rangle  & = E^{0}_n\langle m^{0}|n^{1}\rangle + E^{1}_n\langle m^{0}|n^{0}\rangle \\
    E^{0}_m\langle m^{0}|n^{1}\rangle + \langle m^{0}|H^{1}|n^{0}\rangle & =  E^{0}_n\langle m^{0}|n^{1}\rangle                                    \\
    (E^{0}_n - E^{0}_m) \langle m^{0} | n^{1} \rangle                    & =  \langle m^{0} | H^{1} | n^{0} \rangle                                \\
    \langle m^{0} | n^{1} \rangle                                        & =  \frac{\langle m^{0} | H^{1} | n^{0} \rangle}{E^{0}_n - E^{0}_m}
\end{align*}
Also recall that the unperturbed eigenstates are orthonormal $\braket{m^0  | n^0 }=\delta_{mn}$.
This gives the projection of  $\ket{n^1}$ onto the unperturbed basis vector $\ket{m^0}$, except for the case of $m=n$, which correspond to the parallel of $\ket{n ^{1 }}$ with respect $\ket{m ^{1}}$.

To see this, consider the decomposition of $\ket{n ^{1 }}$ in terms of $\ket{m ^{0}}$
\begin{equation*}
    |n^{1}\rangle = \sum_{m} |m^{0}\rangle \langle m^{0}|n^{1}\rangle
\end{equation*}
However, as we have seen, the projection coefficient is undetermined for $m=n$, hence we split them.
\begin{equation*}
    |n^{1}\rangle = |n^{0}\rangle \langle n^{0}|n^{1}\rangle+\sum_{m \neq n} |m^{0}\rangle \langle m^{0}|n^{1}\rangle
\end{equation*}
Here the first term correspond to the parallel components, as mentioned, while the second to the perpendicular component.

With $\ket{n^1}=\ket{n_\parallel^1}+\ket{n_\perp^1}$ with respect to $\ket{n^0}$ and first order approximation $\ket{n}=\ket{n^0}+\ket{n ^{1}}$, normalization require
\begin{align*}
    \braket{n  | n } & =  \braket{\bra{n^0}+\bra{n_\parallel^1}+\bra{n_\perp^1} | \ket{n^0}+\ket{n_\parallel^1}+\ket{n_\perp^1}}                          \\
    1                & =  \langle n^0 | n^0 \rangle + \langle n^1_{\parallel} | n^0 \rangle + \langle n^0 | n^1_{\parallel} \rangle + \text{higher order}
\end{align*}
Only the orthonormal has real value $\braket{n^0 | n^0 }=1$, so
\begin{equation*}
    0 = \langle n^1 | n^0 \rangle + \langle n^0 | n^1 \rangle + \text{higher order} = 2 \mathrm{Re }\braket{n ^{0 } | n ^{1 }_{\parallel}}+ \text{higher order}
\end{equation*}
This mean  that $\braket{n ^{0 } | n ^{1 }_{\parallel}}= i\alpha$.
By definition, parallel vector mean
\begin{equation*}
    \ket{n_{\parallel} ^{1}} = \braket{n ^{0 } | n ^{1 }} \ket{n ^{0}}
    =\braket{n ^{0 } | n ^{1 }_\parallel} \ket{n ^{0}} + \braket{n ^{0 } | n ^{1 }_\perp} \ket{n ^{0}}
    = i \alpha \ket{n ^{0 }} = i\alpha  \ket{n ^{0 }}
\end{equation*}
where the inner product of $\ket{n ^{1}_\perp}$ is $\ket{n ^{0}}$ is zero.
Substituting this back into the first order approximation
\begin{align*}
    \ket{n}
     & = \ket{n^0}+\ket{n_\parallel^1}+\ket{n_\perp^1}
    =\ket{n^0}+i\alpha\ket{n^0}+\ket{n_\perp^1}
    =(1+i\alpha)\ket{n^0}+\ket{n_\perp^1}                                                                           \\
    \ket{n}
     & =  e ^{i\alpha}\ket{n^0}+\sum_{m \neq n} \frac{ \braket{m^{0} | H^{1} | n^{0} }}{E^{0}_n - E^{0}_m}\ket{m^0}
\end{align*}
By multiplying with $e ^{-i\alpha}$ we can remove the phase factor to the first term, all while the second term being unaffected, since for first order approximation
\begin{equation*}
    (1 - i\alpha)(\text{first order}) = \text{first order} + \text{second order}
\end{equation*}
the second order term $-i\alpha (\text{second order})$  gets  discarded.
Calling the perturbed eigenket with the new phase also $\ket{n }$, we obtain the result to the first order
\begin{equation*}
    \ket{n}
    = \ket{n ^{0 }} + \ket{n ^{1 } }
    = \ket{n _{\parallel}} +\ket{n _{\perp} }
    =\ket{n^0}+\sum_{m \neq n} \frac{ \braket{m^{0} | H^{1} | n^{0} }}{E^{0}_n - E^{0}_m}\ket{m^0}
\end{equation*}

To obtain the second order, we collect the coefficient for  $\epsilon^2$ term
\begin{equation*}
    H^0|n^2\rangle + H^1|n^1\rangle = E^0_n \ket{n^2} + E^1_n \ket{n^1} + E^2_n \ket{n^0}
\end{equation*}
Then dotting with $\ket{n ^{0}}$
\begin{align*}
    E^0_n \braket{n^0 |n^2 }+ \braket{n ^{0 } | H^1|n ^{1}} & =  E_n^{0} \langle n^{0} | n^{2} \rangle + E_n^{1} \langle n^{0} | n^{1} \rangle + E_n^{2} \langle n^{0} | n^{0} \rangle \\
    E_n^{2}                                                 & =      \braket{n ^{0 } | H^1|n ^{1}}                                                                                     \\
                                                            & = \sum_{m \neq n} \frac{\langle n^0 | H | m^0 \rangle \langle m^0 | H | n^0 \rangle}{E^0_n - E^0_m}                      \\
    E_n^{2}
                                                            & = \sum_{m \neq n} \frac{|\langle n^0 | H | m^0 \rangle|^2}{E^0_n - E^0_m}
\end{align*}

\subsection{Symmetry}
Translation operator, which translates the state, defined in abstract state $\ket{\psi}$ and definite position state $\ket{x}$ as 
\begin{equation*}
    T(\epsilon) \ket{\psi} = \ket{\psi_\epsilon}
    T(\epsilon) \ket{x} = \ket{x+\epsilon}
\end{equation*}
That is to say that under the infinitesimal translation $\epsilon$, each state $\ket{\psi}$ gets modified into a translated state $\ket{\psi_\epsilon}$.
The active transformation picture  that corresponds to physically displacing the particle to the right by $\epsilon$ is 
\begin{equation*}
    \langle \psi|T^\dagger(\epsilon) X T(\epsilon)|\psi\rangle = \langle \psi|X|\psi\rangle + \epsilon \qquad
\langle \psi|T^\dagger(\epsilon) P T(\epsilon)|\psi\rangle = \langle \psi|P|\psi\rangle
\end{equation*}
Whereas the passive transformation picture that corresponds to moving the environment to the left by $\epsilon$
\begin{equation*}
    X \to T^{\dagger}(\epsilon)XT(\epsilon) \qquad P \to T^{\dagger}(\epsilon)PT(\epsilon)
\end{equation*}

\begin{table}[]
  \centering
  \caption{Table: Correspondence between Classical and Quantum Mechanical Concepts Related to
Translational lnvariance}
  \begin{tabular}{@{}lll@{}}
    \toprule
    Concept & Classical mechanics & Quantum mechanics \\
    \midrule
    Translation &  $\begin{cases} x\to x+\epsilon\\ p \to p\end{cases}$& $\begin{cases}\braket{X} \to \braket{X}+\epsilon\\\braket{P }\to \braket{P }\end{cases}$\\
    Translational invariance & $H\to H$& $\braket{H }\to \braket{H }$\\
    Conservation law & $\dot{p} = 0$ & $\braket{\dot{P}  }=0$\\
    \bottomrule
  \end{tabular}
\end{table}

We define translational invariance by the requirement
\begin{equation*}
    \langle \psi | H | \psi \rangle = \langle \psi_{\epsilon} | H | \psi_{\epsilon} \rangle 
\end{equation*}
Since $\epsilon$ corresponds to no translation, we may expand $T(\epsilon)$ to order $\epsilon$ as
\begin{equation*}
    T(\epsilon) = I - \frac{i \epsilon}{\hbar} G
\end{equation*}
The operator $G$ is called the generator of transformation.
It is equal to $P$ for spatial translation and equal to $H$ for temporal translation.
It is also equal to $L$ and $S$ for rotation and spin rotation respectively.

The operator $T(a)$ corresponding to a finite translation $a$ can be examined by divide the interval $a$ into $N$ parts of infinite size $a/N$ as $N \to \infty$, and we have $T(a/N) = I -iaP/\hbar N$.
Since a translation by $a$ equals $N$ translations by $a/N$, we find 
\begin{equation*}
    T(a) = \lim_{N \to \infty} \left( I - \frac{iaP }{\hbar N}\right)^N = \exp \left( -iaP/\hbar \right) 
\end{equation*}


\end{document}